

\section{Diagnósticos del Modelo Bayesiano}
\label{app:diagnosticos}

\subsection{Especificación de Priors}

Para el modelo GARCH(1,1), se seleccionaron las siguientes distribuciones a priori (Priors) basándose en la literatura estándar de econometría financiera, buscando distribuciones poco informativas para permitir que los datos dominen la inferencia:

\begin{table}[h]
    \centering
    \begin{tabular}{llp{6cm}}
        \toprule
        \textbf{Parámetro} & \textbf{Distribución A Priori} & \textbf{Justificación} \\
        \midrule
        $\mu$ (Media) & $\mathcal{N}(0, 1)$ & Se asume retorno medio cercano a cero. \\
        $\omega$ (Intercepto) & $\text{InverseGamma}(2.5, 1.0)$ & Positividad estricta requerida. \\
        $\alpha, \beta$ (GARCH) & $\text{Beta}(\dots)$ & Restricción de estacionariedad ($\alpha + \beta < 1$). \\
        $\nu$ (Grados Libertad) & $\text{Gamma}(2, 0.1)$ & Permite colas pesadas (t-Student). \\
        \bottomrule
    \end{tabular}
    \caption{Resumen de la estructura Bayesiana del modelo.}
    \label{tab:priors}
\end{table}

\noindent La elección de estas distribuciones a priori sigue un enfoque \textit{débilmente informativo}, diseñado para imponer restricciones teóricas necesarias sin sesgar excesivamente la inferencia final. A continuación se detalla la lógica detrás de cada selección:

\begin{enumerate}
    \item \textbf{Parámetro de Localización ($\mu$):} Se utiliza una distribución Normal centrada en cero. Dado que estamos modelando variaciones diarias de la prima de riesgo (diferencias), es razonable asumir que el valor esperado a largo plazo tiende a la neutralidad, permitiendo desviaciones simétricas positivas o negativas.
    
    \item \textbf{Restricción de Positividad ($\omega$):} El término de intercepto de la varianza no puede ser negativo. La distribución \textit{InverseGamma} es estándar en la literatura de volatilidad estocástica porque su soporte es estrictamente positivo $(0, \infty)$, evitando así varianzas negativas imposibles físicamente.
    
    \item \textbf{Estacionariedad y Persistencia ($\alpha, \beta$):} Los coeficientes del proceso GARCH representan el peso de las noticias recientes ($\alpha$) y la persistencia de la volatilidad pasada ($\beta$). Ambos deben estar acotados en el intervalo $[0, 1]$. La distribución Beta es la elección natural para variables acotadas en este rango. Además, mediante el uso de \texttt{pm.Potential}, se impuso la restricción dura de estacionariedad $\alpha + \beta < 1$, necesaria para que la varianza no explote al infinito a largo plazo.
    
    \item \textbf{Colas Pesadas ($\nu$):} Los retornos financieros raramente siguen una distribución Normal perfecta (exceso de curtosis). El parámetro de grados de libertad de la t-Student controla el grosor de las colas. Una distribución Gamma para $\nu$ asigna baja probabilidad a valores muy pequeños (varianza infinita) y a valores muy grandes (convergencia a Normal), favoreciendo valores intermedios (típicamente entre 3 y 30) que son característicos de los mercados en crisis.
\end{enumerate}

\section{Registro de Ejecución y Validación Numérica}
\label{app:log_terminal}

Para garantizar la reproducibilidad de los resultados expuestos en la Sección \ref{sec:analisis_empirico}, se adjunta a continuación la salida directa generada por el script de calibración. 

Este registro confirma los tiempos de cómputo (797 segundos), los parámetros exactos calibrados ($\alpha=0.9360$) y detalla las advertencias del muestreador NUTS respecto a la convergencia (R-hat y ESS), justificando la necesidad de mayor potencia computacional discutida en el cuerpo del trabajo.

\begin{lstlisting}[
    basicstyle=\ttfamily\scriptsize, 
    frame=single, 
    backgroundcolor=\color{backcolour}, 
    caption={Salida por terminal del proceso de calibración},
    label={lst:terminal_output}
]
python3 vBayes-student-mejorado.py
-----------------------------------------------------------
--- FASE 1: PREPARACIÓN DE DATOS ---
--> Leyendo germany.xlsx...
    Columna detectada: Bid Yield: Rendimiento bono (oferta) 10 años Germany (%)
--> Leyendo greece.xlsx...
    Columna detectada: Bid Yield: Rendimiento bono (oferta) 10 años Grecia (%)
Datos originales Max: 1987.40 bps
Datos escalados: Max=51.32 | Tipo=float64

--- FASE 2: CALIBRANDO MODELO DE DIFUSIÓN ---
Exponente de difusión Alpha calibrado: 0.9360

--- FASE 3: CALIBRANDO MODELO GARCH (MCMC) ---
Ejecutando calibración Monte Carlo...
Initializing NUTS using adapt_diag...
Multiprocess sampling (2 chains in 2 jobs)
NUTS: [mu, omega, alpha, beta, nu]

Progress: 200 Draws, 0 Divergences, Step size 0.001, Grad evals 167
Sampling Speed: 4.00 s/draws, Elapsed: 0:13:16
Progress: 200 Draws, 0 Divergences, Step size 0.001, Grad evals 191
Sampling Speed: 3.49 s/draws, Elapsed: 0:11:33

Sampling 2 chains for 100 tune and 100 draw iterations (200 + 200 draws total) took 797 seconds.
We recommend running at least 4 chains for robust computation of convergence diagnostics
The rhat statistic is larger than 1.01 for some parameters. This indicates problems during sampling.
The effective sample size per chain is smaller than 100 for some parameters.

Generando simulaciones predictivas (Validación Monte Carlo)...
Sampling: [obs]
Sampling ... 100% 0:00:00 / 0:00:00

===========================================================
      COMPARACIÓN DE MODELOS (MSE AUTOMÁTICO)
===========================================================
MODELO 1: DIFUSIÓN (ECONOPHYSICS)
-> MSE: 8.3865

MODELO 2: GARCH BAYESIANO (MONTE CARLO)
-> MSE: 1692.0925
------------------------------------------------------------
Resultados guardados en analisis_grecia_calibracion_mse
\end{lstlisting}

Enlace al repositorio con el código fuente, datos y más: \url{https://github.com/Leonin04/ModelosComplejosModelosDinamicos}